\documentclass[11pt,a4paper]{article}
\usepackage[margin=25mm]{geometry}
\usepackage{amsmath,amssymb,amsthm,hyperref}
\newtheorem{lemma}{Lemma}
\newtheorem{theorem}[lemma]{Theorem}
\newtheorem{corollary}[lemma]{Corollary}
\newcommand{\gen}[1]{\langle#1\rangle}
\title{Proper face groups in a Suzuki rotation marking}
\author{Human Compression project: a bounded F003 replacement argument}
\date{}
\begin{document}\maketitle

We isolate an obstruction which avoids the intermediate matrix
classification in Spiga's supplied preprint~\cite{Spiga}. It proves the
rank-four exclusion when the field exponent is prime, and identifies
precisely the remaining subfield case. It is not a complete recovery of
that preprint's all-field theorem. The original theorem and its
historical conjecture attribution remain Spiga's; this note grants no
new conjecture or whole-original credit.

Put $S=\operatorname{Sz}(q)$, where $q=2^N$, $N\geq3$ is odd, and
$\theta=2^{(N+1)/2}$. We retain the following classical Suzuki structure,
already used in the project's accepted Suzuki application: the doubly
transitive ovoid action; a unique ovoid fixed point for each involution;
Sylow $2$-groups of exponent four with every involution central and with
each involution centralizer a Sylow $2$-group; and the
maximal-subgroup list consisting of point stabilizers, odd-torus
normalizers and proper Suzuki subfield groups. The torus normalizers
are $D_{2(q-1)}$ and $C_{q\pm\sqrt{2q}+1}\rtimes C_4$. Their odd kernels
are cyclic, and every involution outside that kernel acts on it by
inversion. A nonidentity odd-order element has a unique maximal odd
torus as its centralizer. These are retained structural premises from
Suzuki~\cite{Suzuki}, not classifications reproved here.

In the standard four-dimensional model, the central involutions of the
root group at infinity and the Weyl involution are
\[
 z(s)=\begin{pmatrix}1&0&0&0\\0&1&0&0\\s&0&1&0\\s^\theta&s&0&1\end{pmatrix},
 \qquad J=\begin{pmatrix}0&0&0&1\\0&0&1&0\\0&1&0&0\\1&0&0&0\end{pmatrix}.
\]
The opposite central involutions are $Jz(t)J$. The stabilizer of the
ordered pair $(\infty,0)$ is diagonal, hence has the form
$\operatorname{diag}(d,e,e^{-1},d^{-1})$ with $d,e\ne0$. The Suzuki
relation between $d$ and $e$ will not be needed. These coordinate
descriptions are also retained; the calculation below uses only the
displayed matrices.

\begin{lemma}\label{odd}
The product of two involutions of $S$ is either the identity, an
involution, or an element of odd order.
\end{lemma}
\begin{proof}
The only possible nontrivial even orders in $S$ are two and four, since
an even-order element lies in a Sylow $2$-group or has an odd-order
power which centralizes an involution, and the latter centralizer is a
$2$-group. If involutions $x,y$ had product of order four, then
$\gen{x,y}$ would be a dihedral group of order eight. It lies in a
Sylow $2$-group, where $x$ and $y$ are central and must commute, a
contradiction.
\end{proof}

\begin{lemma}\label{opposite}
Let $s,t,d,e$ be nonzero field elements and put
$D=\operatorname{diag}(d,e,e^{-1},d^{-1})$. Then
$z(s)D Jz(t)J$ is not an involution.
\end{lemma}
\begin{proof}
Direct multiplication gives
\[
 M=z(s)D Jz(t)J=
 \begin{pmatrix}
 d&0&td&t^\theta d\\
 0&e&0&te\\
 sd&0&std+e^{-1}&s t^\theta d\\
 s^\theta d&se&s^\theta td&ste+s^\theta t^\theta d+d^{-1}
 \end{pmatrix}.
\]
The $(1,2)$ entry of $M^2$ is $s t^\theta de$, which is nonzero.
The corresponding entry of the identity is zero.
\end{proof}

\begin{theorem}\label{main}
Suppose $S=\gen{a,b,c}$ and $u=ab$, $v=bc$, $w=abc$ are
involutions. Suppose $|a|,|c|>2$ and $b^2\ne1$. Put
$F=\gen{a,b}$ and $E=\gen{b,c}$. Then $F$ and $E$ cannot both
be contained in subgroups which are either point stabilizers or
odd-torus normalizers.
\end{theorem}
\begin{proof}
The relations give
\[
 a=wv,\qquad c=uw,\qquad F=\gen{a,u},\qquad E=\gen{c,v}.
\]
By Lemma~\ref{odd}, $a,c$ have odd order. If $F$ lies in an
odd-torus normalizer, then $a$ belongs to its odd cyclic kernel and
$u$ inverts $a$. Therefore $b=a^{-1}u$ satisfies $b^2=1$, contrary
to hypothesis. The same argument with $c,v$ excludes a torus
normalizer containing $E$.

Under the contrary assumption, $F$ and $E$ thus lie in point
stabilizers $S_P$ and $S_Q$. If $P=Q$, the three generators belong
to one proper point stabilizer, contrary to generation of $S$.
Conjugate the ordered pair $(P,Q)$ to $(\infty,0)$. The involutions
$u\in F$ and $v\in E$ have these unique fixed points, so
$u=z(s)$ and $v=Jz(t)J$ for nonzero $s,t$. The element
$b\in F\cap E$ belongs to the diagonal two-point stabilizer.
Consequently $w=u b^{-1}v$ has exactly the form excluded by
Lemma~\ref{opposite}. This contradiction proves the theorem.
\end{proof}

\begin{corollary}\label{prime}
If $N$ is an odd prime, $S$ has no rotation marking of a chiral
polytope of rank four.
\end{corollary}
\begin{proof}
We use the standard rotation presentation and chirality criterion:
for a chiral rank-four polytope, $a,b,c$ generate the group,
$ab,bc,abc$ are involutions, the two face subgroups $F,E$ above
are proper, and there is no group automorphism sending
$(a,b,c)$ to $(a^{-1},a^2b,c)$. The last criterion is the usual
extension to the first flag reflection; it is a retained general
polytope input. If $a^2=1$, the identity automorphism would satisfy
it, a contradiction. By duality $c^2\ne1$ as well.

If $b^2=1$, the relations $(ab)^2=(bc)^2=1$ say that $b$
inverts both $a$ and $c$. Then $(abc)^2=1$ reduces to
$ac^{-1}a^{-1}c=1$, so $a$ and $c$ commute. The group
$S$ has the abelian normal subgroup $\gen{a,c}$ and quotient of
order at most two. It is soluble, contradicting its nonabelian
simplicity. Thus all hypotheses of Theorem~\ref{main} hold.

When $N$ is prime, the only possible proper Suzuki subfield group
is $\operatorname{Sz}(2)\cong C_5\rtimes C_4$. Its characteristic
$C_5$ has a unique maximal odd torus as centralizer in $S$;
normalizing $C_5$ therefore normalizes that torus. Hence this
subfield group is already contained in an odd-torus normalizer.
The retained maximal-subgroup list now puts each proper $F,E$ in
one of the two types excluded by Theorem~\ref{main}.
\end{proof}

For general odd $N$, the same argument shows that any putative chiral
rank-four marking must have at least one of its two face groups inside
a proper Suzuki subfield group of parameter at least eight. This is an
explicit remaining obstruction, not a proof that it cannot occur. The
all-rank, almost-simple reduction used by Spiga retains the earlier
Hubard--Leemans theorem; it is not reproved by this note.

The diagnostic programs in the accompanying working directory enumerate
the fixed group $\operatorname{Sz}(8)$. They find no generating rotation
triple with both face groups proper. They motivated the argument, but
are not premises of the written uniform results. No finite search is
used to infer the subfield exclusion still required for composite $N$.

\begin{thebibliography}{3}
\bibitem{Spiga} P. Spiga, \emph{On the rank of Suzuki polytopes: an
answer to Hubard and Leemans}, arXiv:2010.04080v1, 8 October 2020.
The fixed source edition has four pages.
\bibitem{Suzuki} M. Suzuki, \emph{On a class of doubly transitive
groups}, Annals of Mathematics 75 (1962), 105--145.
\bibitem{HL} I. Hubard and D. Leemans, \emph{Chiral polytopes and
Suzuki simple groups}, Rigidity and Symmetry, Fields Institute
Communications 70 (2014), 155--175.
\end{thebibliography}
\end{document}
